\documentclass[12pt]{article}
\usepackage[utf8]{inputenc}
\usepackage[T1]{fontenc}
\usepackage[margin=1in]{geometry}
\usepackage{setspace}
\usepackage{xcolor}
\usepackage[colorlinks=true, urlcolor=blue!60!black, linkcolor=blue!60!black]{hyperref}
\setlength{\parindent}{0em}
\setlength{\parskip}{0.6em}
\renewcommand{\baselinestretch}{1.0}
\begin{document}
\begin{center}
{\LARGE \textit{Palko v.\ Connecticut}, 302 U.S.\ 319 (1937)\par}
\vspace{0.5em}
{\large Case Brief --- Two-Page Edition\par}
\end{center}

\section{Memory Jogger}
Only rights ``implicit in the concept of ordered liberty'' bind the States; Connecticut could appeal and retry a murder defendant (overruled by \textit{Benton}, 1969) (\href{https://supreme.justia.com/cases/federal/us/302/319/}{Justia}).

\section{Detailed Case Facts}
Frank Palka killed two Bridgeport police officers in 1935. Convicted of second-degree murder and sentenced to life, he was retried after the State appealed under an 1886 statute and won reversal for legal errors. The second jury convicted him of first-degree murder, and he was sentenced to death.

\section{Procedural History}
The Connecticut Supreme Court of Errors ordered the retrial and then affirmed the death sentence; Palko appealed as of right.

\section{Judicial Votes}
Affirmed, 8--1. Cardozo wrote for the Court, including Justice Black; Butler dissented without opinion.

\section{Holding}
The Fourteenth Amendment does not incorporate the whole Bill of Rights. It absorbs only rights ``of the very essence of a scheme of ordered liberty''; a state retrial to correct legal error is not ``a hardship so acute and shocking that our polity will not endure it.''

\section{Analysis of Opinions}
Cardozo warned against ``the tyranny of labels.'' Precedent let States drop grand juries, juries, and the self-incrimination privilege, but absorbed free speech and the right to counsel in capital cases. The rationalizing principle is whether liberty and justice could exist without the right; free speech is the ``matrix'' of other freedoms. Connecticut sought only ``a trial free from the corrosion of substantial legal error.'' Butler wrote no opinion.

\section{Examples: Future Applications}
\textbf{Same-Side:} a State charges a felony by prosecutor's information rather than by grand jury indictment. \textbf{Opposite-Side:} a State lets officials enjoin newspapers from criticizing judges. \textbf{Fence-Sitter:} juries that need not be unanimous, allowed under \textit{Palko}'s universal test but barred under the later test asking what is fundamental to the American scheme (\textit{Ramos}, 2020).

\section{Critique}
Total-incorporation originalists (Black, Amar, Thomas) fault its reliance on due process and judicial selection of rights; civil libertarians found its ``civilized system'' benchmark stingy.

\section{Key Quotations}
\begin{quote}
``In these and other situations, immunities that are valid as against the federal government by force of the specific pledges of particular amendments have been found to be implicit in the concept of ordered liberty, and thus, through the Fourteenth Amendment, become valid as against the states.''
\end{quote}
\noindent --- Justice Cardozo, for the Court, 302 U.S.\ at 324--25.

\end{document}
